\section{What the harness actually establishes}
\label{part:harness}

Part~\ref{part:walkthrough} describes a checker that does not exist yet. What
exists is a regression corpus that pins the obligations the checker will have to
discharge. This part explains what that corpus does and --- more usefully --- what
it does not.

\subsection{Two kinds of test}

The suite contains 45 IR fixtures and 7 integration tests, and they split into
two categories that behave completely differently.

\textbf{Green tests are facts locked in.} They assert something true right now
about the IR, the metadata, or the toolchain. They catch regressions.

\textbf{Red tests are executable specifications.} Seventeen fixtures fail on
purpose. Each carries a diagnostic run with an expected diagnostic naming the
exact stable reason a future analysis must emit at a specific operation, and each
fails with \emph{expected error not produced}. They are bring-up gates, not
disabled tests: implementing the analysis and inserting it into those runs is
what turns them green. A suite where every test passes would be a suite that had
stopped specifying anything.

\subsection{Four mechanical layers}

\begin{center}
\begin{tabular}{@{}l p{0.62\linewidth}@{}}
\toprule
Layer & What it establishes \\
\midrule
IR validity     & every fixture parses and satisfies dialect invariants \\
Shape           & the decisive operation --- store, branch, address, division,
                  allocation, backedge --- is present and connected \\
Shape stability & for the eight fixtures carrying a \code{--canonicalize} run
                  --- the four precision controls, the allocation pair, the loop,
                  and the predecessor-choice anti-control --- the decisive shape
                  survives canonicalization, so those cannot silently decay into
                  a different scenario. The other 37 pin shape only in the form
                  they were written \\
Metadata        & each fixture has exactly one of the four outcomes, one
                  observer, a stable reason, obligations required precisely when
                  the outcome is \Unknown{} or \Conditional{}, and real
                  provenance \\
\bottomrule
\end{tabular}
\end{center}

Note what is absent. \textbf{No layer checks an information-flow result},
because no analysis is wired in. The corpus states the specification; it does not
yet test an implementation of it. Being clear about this is the difference between
a useful artifact and a misleading one.

\subsection{Controls, and why every control needs an anti-control}
\label{sec:pairs}

A corpus of leaky programs is satisfied by an analysis that reports every program
as unsafe. That is not a hypothetical failure --- it is the cheapest way to pass a
suite built only from violations, and it makes the tool worthless without making
any test go red.

The guard is a family of \emph{precision controls}: fixtures that are genuinely
noninterferent, where a reported violation is imprecision rather than a finding.
The four in the corpus are exactly the four examples from \cref{step:diag}:
identical-successor branch, value cancellation, public overwrite before
observation, and offset-disjoint reload.

But controls alone are unsafe to ship, because the cheapest way to quiet a
control is usually an unsound repair. So a control wants a paired
\textbf{anti-control} that must retain the opposite outcome:

\begin{center}
\begin{tabular}{@{}p{0.34\linewidth} p{0.54\linewidth}@{}}
\toprule
Control & Unsound repair its partner blocks \\
\midrule
identical-successor branch
  & ``identical successors imply no control leak'' --- the partner is the same
    single-successor branch with differing block-argument operands \\[3pt]
public overwrite
  & adding a proof-authoritative strong update to the diagnostic layer \\[3pt]
offset-disjoint reload
  & coarsening offsets to a cache line, which is a refusal and never a proof \\[3pt]
public world-structural allocation size
  & refusing every dynamic allocation, or treating an equal public cap as an
    equal size \\
\bottomrule
\end{tabular}
\end{center}

The rule is: \emph{never satisfy one member of a pair by weakening the other.}
The first row is the sharpest instance, and it is not a stylistic worry. Running
\code{mlir-opt}'s canonicalizer on an ordinary two-armed diamond produces a
single-successor branch whose edges carry different operands --- byte-identical in
shape to the control, differing only in the operands. An implementation repaired
by that rule accepts a real leak, silently.

\subsection{Which step each fixture guards}

\begin{center}
\begin{tabular}{@{}l p{0.58\linewidth}@{}}
\toprule
Step & Fixtures \\
\midrule
\cref{step:manifest}    & the wrong-party and wrong-host families; the actor
                          examples for the four relations \\
\cref{step:coalitions}  & the downward-closure example, where the leak exists
                          only at a coalition nobody authored \\
\cref{step:normalform}  & currently \emph{uncovered}: no fixture exercises a
                          nonconformance refusal \\
\cref{step:vacuity}     & currently \emph{uncovered}: the vacuity countermodel
                          is not encoded \\
\cref{step:diag}        & the four precision controls \\
\cref{step:product}     & the predecessor-choice anti-control; trace order and
                          multiplicity are \emph{not} yet encoded \\
\cref{step:verdict}     & the unproved-alias and bound-exhaustion refusals \\
\cref{step:deployment}  & the CKKS and wolfSSL helper-latency families, the only
                          \Conditional{} fixtures \\
\bottomrule
\end{tabular}
\end{center}

\subsection{One call, four outcomes}

The wolfSSL 3579 family is the most instructive group in the corpus, because it
demonstrates that the observer model is a load-bearing \emph{input} rather than a
label. \emph{Three} fixtures share one byte-identical body --- a wide multiply
lowered to \fx{@\_\_muldi3} on a 32-bit target --- and differ only in the
declared helper profile. Two more complete the family without sharing that body:
the un-lowered source, whose whole body is a single \code{llvm.mul}, and a
repaired constant-time replacement, which is a 64-iteration mask/add loop that
emits no helper call at all:

\begin{center}
\begin{tabular}{@{}l l p{0.38\linewidth}@{}}
\toprule
Variant & Outcome & Declared knowledge \\
\midrule
source                   & \Unknown{}     & no target lowering, no helper contract \\
target, no contract      & \Unknown{}     & the call exists; no latency summary \\
target, affected profile & \Unsafe{}      & a named operand-dependent helper profile \\
target, constant latency & \Verified{}    & a named constant-latency profile \\
target, repaired         & \Conditional{} & repaired, pending target evidence \\
\bottomrule
\end{tabular}
\end{center}

Three of the four outcome values --- \Unknown{}, \Unsafe{}, \Verified{} --- on
one byte-identical body, driven entirely by declarations. The \Conditional{} row
is different in kind: it comes from a code repair, not a declaration change, and
saying otherwise would overstate an already strong point. That
is the cleanest demonstration available of the point in \cref{step:manifest}:
\Lv{0} is a contract, not proof of the contract, and the same instruction is
\Unknown{}, \Unsafe{}, or \Verified{} according to what has been declared about
the machine it runs on. Note also that the two \Unknown{} rows are genuinely
different scenarios with different missing facts, not one file with an
undecided verdict.

\subsection{Countermodels as regression tests}

Four of the metatheory's required negative results are encoded. These are the
tests a \emph{re-implementation} would fail: they capture reasoning errors rather
than coding errors, and reasoning errors are the ones that survive review.

\begin{center}
\begin{tabular}{@{}p{0.26\linewidth} l p{0.42\linewidth}@{}}
\toprule
Fixture & Outcome & Invalid principle refuted \\
\midrule
predecessor choice     & \Unsafe{}  & an SSA slice is closed around a $\phi$ \\
prefix-causal release  & \Unsafe{}  & a future release may condition an earlier observation \\
alias unproved         & \Unknown{} & unproved ABI separation may be assumed \\
bound exhausted        & \Unknown{} & bounded-run filtering is a sound proof domain \\
\bottomrule
\end{tabular}
\end{center}

The two \Unknown{} rows deserve their outcome explained, because \Unknown{} reads
as a weaker result than it is. Unproved aliasing yields neither safety nor a
counterexample: with no proved disjointness clause and no declared may-alias
realization in the product, asserting \Unsafe{} would itself be unsound, since a
counterexample requires a witness that replays under the exact semantics.
Likewise a reachable bound-exhausted transition that yields no replayed bad
execution is \Unknown{} with the bound-adequacy obligation open --- promoting it
to a counterexample is exactly what \cref{step:verdict} forbids.

The bound-exhaustion fixture is the corpus's only loop with a
\emph{secret-dependent} trip count. Two earlier fixtures already contain
fixed-bound loops with loop-carried block arguments ---
\fx{secret\_embedding\_index.fixed}, a 16-iteration public-induction table scan,
and \fx{wolfssl\_3579\_mul.target\_fixed}, a 64-iteration mask/add multiply ---
so backedge joins were exercised before it. What was absent is a backedge whose
iteration count depends on a secret, and therefore any exercise of bound
adequacy at all.

\subsection{The honest gaps}

The outcome distribution is \Verified{} 21, \Unsafe{} 17, \Unknown{} 5,
\Conditional{} 2. For a tool whose central claim is graceful degradation, five
refusals is thin: there is no fixture for an indirect call, recursion, inline
assembly, a volatile or atomic access, an uninitialized read, or ambiguous
floating-point. Each of those is a place the analysis will meet something it
cannot model, and an untested refusal path is where a silent \Verified{} comes
from.

Three further gaps are worth naming precisely:

\begin{itemize}[nosep,leftmargin=1.4em]
\item About twenty of the older fixtures encode their entire safe/unsafe
      distinction in SSA value names, which \code{mlir-opt} discards at parse.
      After parsing, such a fixture is two indistinguishable stores to two
      unannotated pointers --- so its bring-up gate is not merely unimplemented,
      it is unsatisfiable in principle. The newer fixtures carry policy in
      discardable attributes, which do survive parsing and canonicalization.
\item The corpus is entirely scalar, with three loops, only one of which has a
      non-constant trip count, so nothing exercises
      fixpoint termination at scale, and nothing exercises the vector,
      floating-point, or tensor refusals that any real workload triggers
      immediately.
\item The precision controls carry no diagnostic run, deliberately: their correct
      disposition is \RelReq{}, and a zero-directive run would encode the wrong
      requirement. So they pin shape today and become live precision gates only
      once the record can name an expected disposition. Concretely, the suite
      can currently catch an analysis that \emph{misses} a leak but not one that
      \emph{invents} one.
\end{itemize}
